\section{Problem Formulation}

\subsection{锻造变形序列与状态表征}
铝合金锻造过程中，工件在上、下模具间受压发生塑性变形，其几何形状与温度、等效应力、损伤等物理量随变形历程同步且连续地演化。锻造历程中第 $k$ 步的工件状态以无序点云 $X_k \in \mathbb{R}^{N \times 6}$ 表示，其中 $N$ 为采样点数，每个点由绝对坐标 $p_i=(x_i,y_i,z_i)$ 与温度 $T_i$、等效（von Mises）应力 $\sigma_i$、损伤 $D_i$ 三个物理场组成：
\begin{equation}
X_k = \{\,(p_i,\; T_i,\; \sigma_i,\; D_i)\,\}_{i=1}^{N}
\label{equ:pointcloud}
\end{equation}
其中等效应力由六个应力分量按物理量纲合成：
\begin{equation}
\begin{split}
\sigma = \Big\{ &\tfrac{1}{2}\big[(\sigma_x-\sigma_y)^2+(\sigma_y-\sigma_z)^2 \\
&+(\sigma_z-\sigma_x)^2\big]+3(\tau_{xy}^2+\tau_{yz}^2+\tau_{xz}^2)\Big\}^{1/2}
\end{split}
\end{equation}
温度随热传导与变形生热连续渐变；等效应力与损伤由当前几何与整个变形历程共同决定，具有路径相关性。因此，锻造状态是高维、无序、且几何与多物理场强耦合的点云序列，其演化在时间上连续。

\subsection{控制变量与物理约束}
本文的控制决策变量是整个锻造时间域 $[0,t_f]$（即 $[0,T]$，其中 $T=t_f$ 为锻造总时长）上的完整下压速度函数 $v_y(\cdot)=\{v_y(t)\mid t\in[0,t_f]\}$。该曲线在下压开始前确定，并在锻造过程中开环执行。下压速度通过改变应变速率影响应力场与损伤演化：速度过高引发应变率强化与应力集中，速度过低则延长锻造时间、加剧工件温降与变形抗力。

锻造过程中压机行程 $\Delta h$ 恒定，要求速度曲线在整个锻造时长 $t_f$ 上的积分固定：
\begin{equation}
\int_{0}^{t_f} v_y(t)\,dt = \Delta h = \text{const}
\end{equation}
锻造质量由损伤判据约束。基于 Normalized Cockcroft--Latham 准则，
\begin{equation}
D(t) = \int_0^t \frac{\max(\sigma_1(\tau),0)}{\bar{\sigma}(\tau)}\,\dot{\bar{\varepsilon}}(\tau)\,d\tau
\end{equation}
记临界损伤值为 $D_{\text{cr}}$，损伤累积值一旦超过它即视为开裂。锻造要求全程损伤不超该临界值，即 $D(t) < D_{\text{cr}},\ \forall t\in[0,t_f]$。

在上述物理过程与约束下，本文要解决的问题可分解为如下两个相互关联的子问题。

\textbf{问题一（连续演化建模）}：给定初始步点云 $X_0$ 与下压速度曲线 $v_y(t)$，如何预测工件几何与温度、等效应力、损伤等多物理场在变形历程中的连续演化，并获得紧凑、可推演的状态表征，使模型既能由此预测未来各步的状态，又能高保真地还原为几何—多物理场点云？
    
\textbf{问题二（连续速度曲线决策）}：在恒定锻造行程与损伤约束下，如何决策整条连续速度曲线 $v_y(t)$，使锻造时间 $t_f$ 最短，同时保证损伤全程不超过临界值 $D_{\text{cr}}$？

综合两个子问题，最终要求解的下压速度决策问题可写成如下约束优化形式：
\begin{equation}
\begin{gathered}
\min\nolimits_{v_y(t)} \; t_f \\
\text{s.t.} \quad \int_{0}^{t_f} v_y(t)\,dt = \Delta h \\
D(t) < D_{\text{cr}}, \quad \forall t \in [0, t_f]
\end{gathered}
\label{equ:model}
\end{equation}
即在恒定锻造行程前提下，兼顾生产效率与成形质量，寻求最优速度曲线。
